\clearpage
\begin{center}
    {\headingfont\LARGE Abstract}

\end{center}
\vspace{1.5cm}
\par
A blind person in Bangladesh who wants to read a label, or to know which banknote is in the hand and whether it is real, usually has to ask someone. The assistive devices that could help are expensive, built for English, and often need a phone or the internet. This thesis presents \textit{Savior Glass}, a Bangla-first assistive smart glass built on a Raspberry Pi~5 with a camera, five push-buttons and audio output. It has five modes: Bangla and English text reading, object announcement, Bangladeshi Taka recognition with a cautious counterfeit check, facial-expression announcement, and an optional online scene description. Everything except the last runs on the device.

To recognise notes, we pasted 5,073 background-free note images onto real photographs and let the program write the labels, which gave 22,333 training images without any manual annotation. A YOLOv8s detector trained on them reaches 0.995 mAP@0.5 and 0.849 mAP@0.5:0.95 on 2,282 held-out composites. On 1,536 photographs taken by other people it names the right denomination 91.5\% of the time, but on hand-held close-ups only 18.5\%.

For counterfeit notes we worked with the JaalTaka dataset (1,390 notes, six views each). A model trained on all six views breaks down when a user captures fewer; training on random prefixes of the views fixes this and raises one-view accuracy from 71.2\% to 98.2\% over three seeds. We then found that the usual split of this dataset lets the same counterfeit print appear in both training and test data. On a split with prints the model has never seen, accuracy falls to about 90\%. Checking the banknote watermark, found by a learned 4.3~MB localizer, brings it back to 95.0\%, level with a fine-tuned ResNet-50. Because a wrong verdict can hurt a blind user, the glass never says ``counterfeit''. It says ``likely genuine'' only above a threshold fixed on validation data and otherwise asks for a check by hand; none of 88 test counterfeits was passed as genuine.

For text reading we built a benchmark with correctly shaped Bangla and with separate phrases and fonts for validation and test. Keeping only the main block of text, tuning the detector and cleaning the output lowers the character error rate from 9.3\% to 2.4\% overall and from 11.9\% to 3.2\% for Bangla.

On the Raspberry Pi~5 the currency check takes a median of 327~ms and the system ran for 30~minutes without thermal throttling, but reading a text takes about 19~seconds. In the first hands-on test with the glass camera, currency recognition was unreliable and Bangla text was not read, so the laboratory figures do not yet describe the device as worn. A study with visually impaired participants is prepared, including a voice-driven questionnaire; trial sessions have started, but their results are not yet reported. Every number in this thesis comes from a saved result file, and results that went against us are reported.

\vspace{0.5cm}
\noindent\textbf{Keywords:} Assistive Technology, Smart Glass, Visually Impaired, Raspberry Pi, Bangla OCR, YOLOv8, Bangladeshi Currency, Counterfeit Detection, Multi-View Learning, Watermark, Selective Prediction, Edge AI
\clearpage
